% Prozentrechnung · Veränderung · Prüfungs-Fokus MSA – bank/prozentrechnung/e5.jsonl (zusammenbau v0.6, Prüfungs-Fokus)
\einheitenkopf[e5]{Einheit 5 · Prozentuale Veränderung}
\zweigzeile{ab Kl. 7, je nach Buch bis Kl. 8 · P10 ×7}
\setcounter{aufgabe}{0}

% Anlauf (ohne Prüfkennung)
\begin{aufgabe}{Veränderung – Anlauf}
\begin{teile}
\teil um $20\,\%$ gesenkt – was bedeutet dasselbe? Kreuze an.\\ \kreuz{auf $20\,\%$ gesenkt}\\ \kreuz{auf $80\,\%$ gesenkt}\\ \kreuz{um $80\,\%$ gesenkt}
\teil $80\,€$ um $10\,\%$ erhöht – neuer Preis? \leerfeld[€]
\teil $45\,€$ um $20\,\%$ erhöht – mit Faktor \leerfeld[€]
\end{teile}
\end{aufgabe}

\begin{aufgabe}{Veränderung – Prüfungsaufgaben}
\begin{teile}
\teil Ein Liter Milch kostete $1{,}15\,€$, jetzt kostet er $1{,}29\,€$. Um wie viel Prozent ist der Preis gestiegen? Runde auf eine Stelle. \hfill (P26F)
\teil Eine Kugel Eis kostete $1{,}70\,€$, jetzt kostet sie $1{,}90\,€$. Um wie viel Prozent ist der Preis gestiegen? Runde auf eine Stelle. \hfill (P26F)
\teil Ein Kino hatte im Mai $1\,340$ Besucher, im Juni $1\,120$. Um wie viel Prozent ist die Zahl gesunken? Runde auf eine Stelle. \hfill (P22)
\teil Eine Bücherei verlieh im März $860$ Bücher, im April $780$. Um wie viel Prozent ist die Zahl gesunken? Runde auf eine Stelle. \hfill (P22)
\teil Ein Online-Shop hatte $48\,000$ Bestellungen, im Jahr darauf $57\,600$. Um wie viel Prozent ist die Zahl gestiegen? \hfill (P16)
\teil Eine Stadt hatte $64\,000$ Einwohner, zehn Jahre später $72\,000$. Um wie viel Prozent ist die Zahl gestiegen? \hfill (P16)
\teil Ein Busticket kostet $2{,}40\,€$ und wird um $25\,\%$ teurer. Wie viel kostet es dann? \hfill (P24) \\ \leerfeld[€]
\teil Ein Schwimmkurs kostet $65\,€$ und wird um $12\,\%$ teurer. Wie viel kostet er dann? \hfill (P24) \\ \leerfeld[€]
\teil Kletterhalle: Erwachsene $14\,€$, Jugendliche (bis $17$ Jahre) $9\,€$. Familie Demir: Mutter, Vater, Tochter ($15$ Jahre), Sohn ($12$ Jahre). Montags gibt es $15\,\%$ Rabatt. Wie viel zahlt die Familie am Montag? \hfill (P15)
\teil Theater: Erwachsene $22\,€$, Schüler $13\,€$. Es kommen Oma, Mutter, Vater und ein Schüler. Mit der Familienkarte gibt es $25\,\%$ Rabatt. Wie viel zahlen sie zusammen? \hfill (P15)
\teil Für Radwege gilt: Steigung höchstens $5\,\%$. Ergänze: „$5\,\%$ Steigung heißt: auf \leerfeld[m] waagerechter Strecke \leerfeld[m] Höhenunterschied.“ Herr Lenz sagt: „Ein Radweg, der auf $240$ m waagerechter Strecke um $15$ m ansteigt, ist zu steil.“ Hat er recht? Rechne. \hfill (P25)
\teil Eine Garagenzufahrt darf höchstens $15\,\%$ Steigung haben. Ergänze: „$15\,\%$ Steigung heißt: auf $100$ m waagerechter Strecke \leerfeld[m] Höhenunterschied.“ Eine Zufahrt überwindet auf $8$ m waagerechter Strecke $1{,}1$ m. Frau Roth sagt, sie sei zu steil. Hat sie recht? Rechne. \hfill (P25)
\end{teile}
\end{aufgabe}

\medskip
\uebersichtskasten{Erhöhung um p \%: neuer Wert = alter Wert · (1 + p/100); Senkung: · (1 − p/100). \\ 250 € um 12 \% erhöht: 250 · 1,12 = 280 € \quad 250 € um 12 \% gesenkt: 250 · 0,88 = 220 € \\ Veränderung in Prozent: (neu − alt) : alt. \quad von 60 auf 69: 9 : 60 = 0,15 = 15 \% \\ Prozentpunkte: Unterschied zweier Prozentangaben. \quad von 30 \% auf 33 \%: 3 Prozentpunkte (das sind 10 \% mehr) \\ Steigung in Prozent: Höhe : waagerechte Strecke. \quad 8 m auf 100 m = 8 \%}
